\documentclass[11pt,a4paper]{article}

\usepackage[margin=1.05in]{geometry}
\usepackage{amsmath,amssymb,amsthm}
\usepackage{mathtools}
\usepackage{microtype}
\setlength{\emergencystretch}{2.5em}
\usepackage{enumitem}
\usepackage{array}
\usepackage{booktabs}
\usepackage{tikz-cd}
\usepackage[colorlinks=true,linkcolor=blue!45!black,urlcolor=blue!60!black]{hyperref}

\newtheorem{theorem}{Theorem}[section]
\newtheorem{proposition}[theorem]{Proposition}
\newtheorem{lemma}[theorem]{Lemma}
\newtheorem{corollary}[theorem]{Corollary}
\theoremstyle{definition}
\newtheorem{definition}[theorem]{Definition}
\newtheorem{example}[theorem]{Example}
\newtheorem{notation}[theorem]{Notation}
\theoremstyle{remark}
\newtheorem{remark}[theorem]{Remark}
\newtheorem{warning}[theorem]{Warning}

\newcommand{\FF}{\mathbb{F}}
\newcommand{\CC}{\mathbb{C}}
\newcommand{\RR}{\mathbb{R}}
\newcommand{\QQ}{\mathbb{Q}}
\newcommand{\NN}{\mathbb{N}}
\newcommand{\ZZ}{\mathbb{Z}}
\newcommand{\HH}{\mathbb{H}}
\newcommand{\ip}[2]{\langle #1, #2 \rangle}
\newcommand{\dual}[1]{#1^{*}}
\newcommand{\perpp}{^{\perp}}
\newcommand{\cat}[1]{\mathbf{#1}}
\newcommand{\Set}{\cat{Set}}
\newcommand{\FinSet}{\cat{FinSet}}
\newcommand{\Bool}{\cat{Bool}}
\newcommand{\FinBool}{\cat{FinBool}}
\newcommand{\CABA}{\cat{CABA}}
\newcommand{\Vect}{\cat{Vect}}
\newcommand{\FinVect}{\cat{FinVect}}
\newcommand{\Stone}{\cat{Stone}}
\newcommand{\LinComp}{\cat{LinComp}}
\newcommand{\PG}{\mathrm{PG}}
\newcommand{\op}{\mathrm{op}}
\DeclareMathOperator{\disc}{disc}
\DeclareMathOperator{\Gal}{Gal}
\DeclareMathOperator{\Id}{Id}
\DeclareMathOperator{\Hom}{Hom}
\DeclareMathOperator{\Ob}{Ob}
\DeclareMathOperator{\Aut}{Aut}
\DeclareMathOperator{\Sub}{Sub}
\DeclareMathOperator{\Ind}{Ind}
\DeclareMathOperator{\Pro}{Pro}
\DeclareMathOperator{\Spec}{Spec}
\DeclareMathOperator{\PGL}{PGL}
\DeclareMathOperator{\PGammaL}{P\Gamma L}
\DeclareMathOperator{\GL}{GL}
\DeclareMathOperator{\Sp}{Sp}
\newcommand{\qbin}[2]{\genfrac{[}{]}{0pt}{}{#1}{#2}_{q}}

\title{\bfseries What Replaces the Inner Product over a Finite Field\\[0.4em]
\large and what the opposite category of linear algebra is}
\author{Robin Langer \and Claude}
\date{11 September 2026}

\begin{document}
\maketitle

\begin{abstract}
Part~I asks what survives of the inner product when the ground field is finite. Separating
\emph{nondegeneracy} from \emph{positive-definiteness}, we find the first survives in
abundance and the second cannot even be stated; what is lost is the orthogonal splitting
$V = W \perp W\perpp$, and what is gained is isotropy, hence self-dual codes, polar
geometries and the finite classical groups.

Part~II asks a categorical question with the same shape. The opposite of the category of
sets is the category of complete atomic Boolean algebras. What, then, is the opposite of
the category of vector spaces over $\FF_q$? The answer forks: as a bare category
$\FinVect_q$ is self-dual, but if one deforms the \emph{recipe} behind $\Set^{\op} \simeq
\CABA$ --- take subobject lattices --- one lands on projective geometry, whose
anti-automorphisms are classified by exactly the three families of forms found in Part~I.
The two halves of the document turn out to be one theorem seen twice --- and, by
Piron's and Sol\`er's reconstruction theorems, a third time, as the list of possible quantum
mechanics, the symplectic case being excluded by precisely the definiteness Part~I is about
losing.

Part~III takes $q \to 1$. Isotropy vanishes, the polarity $W \mapsto W\perpp$ becomes
Boolean complementation, projective geometry becomes the Boolean lattice, and we recover
$\Set^{\op} \simeq \CABA$. We also separate the field with one element from the prime at
infinity, which are commonly and productively confused.

\medskip
\noindent\emph{Everything is defined from scratch. A reader who knows what a set is can
read this straight through.}
\end{abstract}

\tableofcontents

\newpage
\part{Forms over a finite field}

\section{The mishnah: what ``naturally isomorphic'' is doing here}

\subsection{The starting objects}

\begin{definition}[Field]
A \emph{field} $k$ is a set with two operations $+$ and $\cdot$, both associative and
commutative, with $\cdot$ distributing over $+$, with additive identity $0$ and
multiplicative identity $1 \neq 0$, such that every element has an additive inverse and
every nonzero element has a multiplicative inverse. Examples: $\QQ$, $\RR$, $\CC$, and for
each prime $p$ the set $\FF_p = \ZZ/p\ZZ$ of integers modulo $p$.
\end{definition}

\begin{definition}[Finite field]
For each prime power $q = p^{e}$ there is, up to isomorphism, exactly one field with $q$
elements, written $\FF_q$. Its \emph{characteristic} is $p$, meaning $\underbrace{1 + \cdots
+ 1}_{p} = 0$. Its multiplicative group $\FF_q^{\times} = \FF_q \setminus \{0\}$ is cyclic
of order $q - 1$. No field has a number of elements that is not a prime power.
\end{definition}

\begin{definition}[Vector space, linear map, dual]
A \emph{vector space} over $k$ is an abelian group $V$ with a scalar multiplication
$k \times V \to V$ satisfying the usual axioms. A \emph{linear map} $f : V \to W$ satisfies
$f(u+v) = f(u)+f(v)$ and $f(\lambda v) = \lambda f(v)$. The \emph{dual space} is
\[
  \dual{V} \;=\; \Hom_k(V, k) \;=\; \{\text{linear maps } V \to k\},
\]
itself a vector space. If $\dim V = n < \infty$ then $\dim \dual{V} = n$, so
$V \cong \dual{V}$ --- but only after choosing a basis.
\end{definition}

\begin{definition}[Bilinear form]
A \emph{bilinear form} on a $k$-vector space $V$ is a map $\ip{-}{-} : V \times V \to k$
linear in each argument separately. It is \emph{symmetric} if $\ip{v}{w} = \ip{w}{v}$,
\emph{alternating} if $\ip{v}{v} = 0$ for all $v$, and \emph{nondegenerate} if
$\ip{v}{w} = 0$ for all $w$ forces $v = 0$.
\end{definition}

\begin{definition}[Gram matrix]
Given a basis $e_1, \dots, e_n$ of $V$, the \emph{Gram matrix} is
$G = (\ip{e_i}{e_j})_{i,j}$. The form is nondegenerate iff $\det G \neq 0$. Changing basis
by an invertible $P$ replaces $G$ by $P^{\mathsf{T}} G P$, so $\det G$ changes by a nonzero
square.
\end{definition}

\subsection{The two ingredients of an inner product}

Let $V$ be a finite-dimensional complex inner product space. The form gives a map
\[
  \flat : V \longrightarrow \dual{V}, \qquad v \longmapsto \ip{-}{v},
\]
conjugate-linear (or linear into the conjugate dual $\overline{V}^{*}$, depending on
convention), and it is an isomorphism because the form is \emph{nondegenerate}.
Positive-definiteness --- the requirement $\ip{v}{v} > 0$ for $v \neq 0$ --- is sufficient
for nondegeneracy but by no means necessary. So the isomorphism $V \cong \dual{V}$ consumes
only a fraction of what an inner product supplies.

Two remarks on the word ``natural'', both of which sharpen the question. (Naturality is
defined precisely in \S\ref{sec:cats}; the reader may take the following on trust for now
and return.)

\begin{remark}
There is no natural transformation $\Id \Rightarrow (-)^{*}$ on finite-dimensional vector
spaces: $(-)^{*}$ is contravariant, so the statement does not typecheck as naturality. What
one has is \emph{canonicity relative to a chosen form}, which has exactly the same epistemic
status as a chosen basis. The genuinely natural isomorphism is $V \cong V^{**}$.
\end{remark}

\begin{remark}\label{rem:ordering}
Consequently the right question is not ``can we define an inner product over $\FF_q$?'' but:
\emph{which of the two ingredients survives?} Nondegeneracy is pure linear algebra and
survives. Definiteness is a statement about an \emph{ordering} of the ground field, and
$\FF_q$ admits none compatible with its arithmetic: in an ordered field $1 > 0$ forces
$1 + 1 + \cdots > 0$, whereas in $\FF_p$ the sum of $p$ copies of $1$ is $0$.
\end{remark}

\section{What can be defined over $\FF_q$}\label{sec:forms}

Three families of nondegenerate forms exist, and --- unlike over $\RR$ or $\CC$ --- all
three are classified completely. Finite fields are, in this respect, \emph{nicer} than the
classical ground fields.

\subsection{Symmetric bilinear and quadratic forms}

\begin{definition}[Isometry]
Two forms $(V, \ip{-}{-}_V)$ and $(W, \ip{-}{-}_W)$ are \emph{isometric} if there is a linear
isomorphism $f : V \to W$ with $\ip{f(u)}{f(v)}_W = \ip{u}{v}_V$. The \emph{isometry group}
of a form is the group of self-isometries.
\end{definition}

\begin{definition}[Discriminant]
For a nondegenerate symmetric form with Gram matrix $G$, the \emph{discriminant} is the class
of $\det G$ in $k^{\times}/(k^{\times})^{2}$. By the change-of-basis rule above this is
independent of the basis.
\end{definition}

For $q$ odd, $\FF_q^{\times}$ is cyclic of even order $q-1$, so squares have index $2$ and
$\FF_q^{\times}/(\FF_q^{\times})^{2} \cong \ZZ/2$.

\begin{theorem}[Classification, $q$ odd]
Nondegenerate symmetric bilinear forms on $\FF_q^{n}$ are classified up to isometry by the
discriminant, an element of $\FF_q^{\times}/(\FF_q^{\times})^{2} \cong \ZZ/2$. There are
exactly \emph{two} isometry classes in each dimension.
\end{theorem}

Compare: over $\CC$ there is exactly one class in each dimension (Gram matrix $I$); over
$\RR$, Sylvester's law of inertia gives $n+1$ classes indexed by signature. Finite fields sit
strictly between, retaining a vestigial ghost of signature of size~$2$. The isometry groups
are the orthogonal groups $O_n^{+}(q)$ and $O_n^{-}(q)$.

\begin{remark}[Characteristic two]
In characteristic $2$ the symmetric and alternating conditions collide, since $x = -x$. The
correct object is then the \emph{quadratic form} $Q : V \to k$ (satisfying $Q(\lambda v) =
\lambda^{2}Q(v)$ and with $B(x,y) = Q(x+y) - Q(x) - Q(y)$ bilinear) rather than its
\emph{polarisation} $B$, which is alternating and carries strictly less information. Read
everything below about symmetric forms with $q$ odd unless stated otherwise.
\end{remark}

\subsection{Alternating (symplectic) forms}

Nondegenerate alternating forms exist only in even dimension, and there is exactly \emph{one}
up to isometry: an orthogonal sum of hyperbolic planes (Definition~\ref{def:hyp}). The
isometry group is the symplectic group $\Sp_{2n}(q)$. The classification is even coarser than
over $\CC$: symplectic geometry has no local invariants at all.

\subsection{Hermitian forms: the honest analogue}

This is the closest finite-field relative of the complex inner product.

\begin{definition}[Galois group, Frobenius]
For a field extension $L \supseteq k$, the \emph{Galois group} $\Gal(L/k)$ is the group of
field automorphisms of $L$ fixing $k$ pointwise. For finite fields, $\Gal(\FF_{q^{d}}/\FF_q)$
is cyclic of order $d$, generated by the \emph{Frobenius} $\sigma : x \mapsto x^{q}$.
\end{definition}

Complex conjugation is the nontrivial element of $\Gal(\CC/\RR)$, a group of order two. The
finite analogue is immediate: take $\FF_{q^{2}} \supset \FF_q$, whose Galois group is
generated by $\sigma : x \mapsto x^{q}$, an involution ($\sigma^{2} = \mathrm{id}$ since
$x^{q^{2}} = x$ on $\FF_{q^{2}}$) with fixed field exactly $\FF_q$ (the roots of $x^{q} = x$).

\begin{definition}[Semilinear, sesquilinear, Hermitian]\label{def:sesq}
Let $\sigma$ be a field automorphism of $L$. A map $f$ is \emph{$\sigma$-semilinear} if
$f(u+v) = f(u)+f(v)$ and $f(\lambda v) = \sigma(\lambda) f(v)$. A
\emph{$\sigma$-sesquilinear form} is $\ip{-}{-} : V \times V \to L$ linear in the first
argument and $\sigma$-semilinear in the second. It is \emph{Hermitian} if $\sigma$ is an
involution and $\ip{w}{v} = \sigma\ip{v}{w}$.
\end{definition}

\begin{proposition}
For a Hermitian form over $\FF_{q^{2}}$ and any $v \in V$ we have $\ip{v}{v}^{q} = \ip{v}{v}$,
hence $\ip{v}{v} \in \FF_q$.
\end{proposition}

That is: \emph{the norm is always real}, exactly as over $\CC$. Moreover there is exactly one
nondegenerate Hermitian form per dimension up to isometry, with Gram matrix $I$; the isometry
group is the unitary group $U_n(q)$. Every structural feature of the complex inner product is
reproduced save one: $\FF_q$ is not ordered, so ``$\ip{v}{v} > 0$'' is meaningless.

Together with $\GL_n(q)$, these isometry groups are the \emph{finite classical groups}. The
question ``what replaces the inner product over $\FF_q$?'' is therefore an entrance to the
classification of finite simple groups.

\section{The objection: what actually breaks}\label{sec:breaks}

Nondegeneracy buys $V \cong \dual{V}$. It does \emph{not} buy the orthogonal splitting.

\begin{definition}[Orthogonal complement, isotropy]
For a subspace $W \subseteq V$ put $W\perpp = \{v \in V : \ip{v}{w} = 0 \ \forall w \in W\}$.
A vector $v \neq 0$ is \emph{isotropic} if $\ip{v}{v} = 0$; a subspace is \emph{totally
isotropic} if the form vanishes identically on it; a form is \emph{anisotropic} if it has no
isotropic vectors.
\end{definition}

If $v$ is isotropic then $v \in \langle v\rangle \cap \langle v\rangle\perpp$, so $V \neq
\langle v\rangle \oplus \langle v\rangle\perpp$. Gram--Schmidt fails: one cannot normalise a
vector of norm zero, nor divide by it. Orthogonal projection, least squares and the analytic
spectral theorem fail with it. There is no norm, no metric, no completeness, no analysis.

Worse --- or better --- isotropy is unavoidable.

\begin{theorem}[Chevalley--Warning, quadratic case]
Every quadratic form over $\FF_q$ in $n \geq 3$ variables has a nontrivial zero. Consequently
an anisotropic quadratic space over $\FF_q$ has dimension at most $2$.
\end{theorem}

Over $\RR$, anisotropic spaces exist in every dimension --- that is precisely what
positive-definiteness means. Over $\FF_q$ they are confined to dimension $\leq 2$.

What does survive, from nondegeneracy alone, is the numerical duality:
\[
  \dim W + \dim W\perpp = \dim V, \qquad (W\perpp)\perpp = W .
\]
\emph{One keeps the duality and loses the splitting.} This sentence is the whole of Part~I in
compressed form, and Part~II is its categorical restatement.

\section{Attempted resolution: Witt's theory}

\begin{definition}[Hyperbolic plane]\label{def:hyp}
The \emph{hyperbolic plane} $\HH$ has basis $e, f$ with $\ip{e}{e} = \ip{f}{f} = 0$ and
$\ip{e}{f} = 1$. It is nondegenerate and contains isotropic vectors by construction.
\end{definition}

\begin{theorem}[Witt decomposition]
Every nondegenerate space decomposes as
\[
  V \;\cong\; \underbrace{\HH \perp \cdots \perp \HH}_{m} \;\perp\; V_{\mathrm{an}},
\]
with $V_{\mathrm{an}}$ anisotropic and unique up to isometry. The integer $m$ is the
\emph{Witt index}, the common dimension of the maximal totally isotropic subspaces.
\end{theorem}

\begin{theorem}[Witt extension]
Any isometry between subspaces of a nondegenerate space extends to an isometry of the whole
space.
\end{theorem}

Witt extension replaces the statement that the unitary group acts transitively on orthonormal
frames: the isometry group acts transitively on isometric configurations. Over $\FF_q$, since
$\dim V_{\mathrm{an}} \leq 2$, finite quadratic spaces are \emph{almost entirely hyperbolic}.
A situation that from the $\CC$ side looks maximally degenerate is in fact the source of all
the interesting combinatorics.

\section{\emph{Mai nafka minah}: what changes in practice}

\subsection{Coding theory and self-duality}

\begin{definition}[Linear code]
A \emph{linear code} of length $n$ over $\FF_q$ is a subspace $C \subseteq \FF_q^{n}$. Its
\emph{dual} is $C\perpp$ taken with respect to the standard bilinear form $\sum_i x_i y_i$.
\end{definition}

\emph{Self-dual} codes, $C = C\perpp$, exist --- the extended binary Golay code $[24,12,8]$,
the hexacode over $\FF_4$ --- and they are the most beautiful objects in the subject. Over
$\RR$ with a positive-definite form this is flatly impossible: $W = W\perpp$ forces $W = 0$.
Self-duality is a gift of isotropy.

\begin{warning}
Self-duality is not special to finite fields; it is special to \emph{indefinite or
unordered} ground fields. Over $\CC$ with the \emph{bilinear} (not Hermitian) form
$x_1y_1 + x_2y_2$, the line $\langle (1,i)\rangle \subset \CC^{2}$ satisfies $W = W\perpp$.
What is special about $\FF_q$ is that Chevalley--Warning makes isotropy \emph{unavoidable}
and abundant.
\end{warning}

\subsection{Finite geometry}

The totally isotropic subspaces of a nondegenerate form, ordered by inclusion, form a
\emph{polar space}: generalised quadrangles and their higher analogues, the geometries on
which the classical groups act. Under positive-definiteness this geometry is empty.

\subsection{Quantum error correction}

Stabiliser codes are totally isotropic subspaces of $\FF_2^{2n}$ under the standard symplectic
form. The requirement that stabiliser generators commute is \emph{exactly} the isotropy
condition, and the quotient $C\perpp/C$ carries the logical operators. Symplectic
$\FF_2$-geometry is the native language of the Clifford group.

\section{Where the analysis actually lives}\label{sec:fourier}

If one wants Fourier analysis, the productive move is to stop trying to place an inner product
on $\FF_q^{n}$ and place it instead on the $\CC$-vector space of functions:
\[
  \ip{f}{g} \;=\; \sum_{x \in \FF_q^{n}} f(x)\,\overline{g(x)},
  \qquad f, g : \FF_q^{n} \to \CC .
\]
This is a genuine positive-definite Hermitian inner product.

\begin{definition}[Additive character]
An \emph{additive character} of $\FF_q$ is a group homomorphism $\psi : (\FF_q, +) \to
(\CC^{\times}, \cdot)$. For $\FF_p$ one may take $\psi(x) = e^{2\pi i x/p}$; in general
$\psi(x) = e^{2\pi i \,\mathrm{Tr}(x)/p}$ with $\mathrm{Tr} : \FF_q \to \FF_p$ the trace.
\end{definition}

Fixing a nontrivial $\psi$, the functions $e_y(x) = \psi\!\left(\sum_i x_i y_i\right)$, for
$y \in \FF_q^{n}$, form an orthogonal basis, giving a Fourier transform, Plancherel's theorem
and Poisson summation.

Observe the division of labour: the $\FF_q$-bilinear form appears \emph{inside the character},
indexing the basis, rather than defining the geometry of the space. The MacWilliams identity,
relating the weight enumerator of a code $C$ to that of $C\perpp$, is Poisson summation in
exactly this setting --- the bilinear form supplies the duality $C \leftrightarrow C\perpp$,
and the $\CC$-valued inner product supplies the analysis.

\section{Summary of Part I}

\begin{center}
\begin{tabular}{lccc}
\toprule
 & $\CC$ & $\RR$ & $\FF_q$ ($q$ odd) \\
\midrule
Nondegenerate symmetric forms & $1$ class & $n+1$ classes & $2$ classes \\
Nondegenerate alternating forms & $1$ class & $1$ class & $1$ class \\
Nondegenerate Hermitian forms & $n+1$ classes & --- & $1$ class \\
$V \cong \dual{V}$ & yes & yes & yes \\
$V = W \perp W\perpp$ & yes & yes & \textbf{no} \\
Anisotropic spaces & $\dim \leq 1$ & all $n$ & $\dim \leq 2$ \\
Self-dual subspaces $W = W\perpp$ & possible & no (definite) & \textbf{yes} \\
Norm, metric, analysis on $V$ & yes & yes & no \\
\bottomrule
\end{tabular}
\end{center}

\newpage
\part{Categories, lattices, and two kinds of opposite}

\noindent
Part~I answered a question about forms. Part~II asks a question of the same shape one level
up, and the answer will turn out to be the same theorem.

\medskip
\noindent\textbf{The question.} The opposite of the category of sets is the category of
complete atomic Boolean algebras. What is the opposite of the category of vector spaces over
$\FF_q$?

\section{Categories, from scratch}\label{sec:cats}

\begin{definition}[Category]
A \emph{category} $\mathcal{C}$ consists of:
\begin{itemize}[nosep]
  \item a collection $\Ob(\mathcal{C})$ of \emph{objects};
  \item for each pair $X, Y$ of objects a set $\mathcal{C}(X,Y)$ of \emph{morphisms}
        (written $f : X \to Y$);
  \item for each $X$ an \emph{identity} $\mathrm{id}_X \in \mathcal{C}(X,X)$;
  \item a \emph{composition} $\mathcal{C}(Y,Z) \times \mathcal{C}(X,Y) \to \mathcal{C}(X,Z)$,
        $(g,f) \mapsto g \circ f$;
\end{itemize}
such that composition is associative and the identities are neutral:
$h(gf) = (hg)f$, $f \circ \mathrm{id}_X = f = \mathrm{id}_Y \circ f$.
\end{definition}

\begin{example}\label{ex:cats}
\begin{enumerate}[nosep,label=(\alph*)]
  \item $\Set$: objects are sets, morphisms are functions. $\FinSet$: the finite ones.
  \item $\Vect_k$: objects are $k$-vector spaces, morphisms are linear maps.
        $\FinVect_k$: the finite-dimensional ones.
  \item Any partially ordered set $(P,\leq)$ is a category: objects are the elements, and
        there is exactly one morphism $x \to y$ when $x \leq y$, none otherwise.
        Transitivity gives composition, reflexivity the identities.
  \item Any monoid is a category with one object, its elements being the morphisms.
\end{enumerate}
\end{example}

\begin{definition}[Isomorphism]
$f : X \to Y$ is an \emph{isomorphism} if there is $g : Y \to X$ with $gf = \mathrm{id}_X$
and $fg = \mathrm{id}_Y$. In $\Set$ these are the bijections; in $\Vect_k$, the invertible
linear maps.
\end{definition}

\begin{definition}[Opposite category]\label{def:op}
The \emph{opposite} $\mathcal{C}^{\op}$ has the same objects as $\mathcal{C}$ and
$\mathcal{C}^{\op}(X,Y) = \mathcal{C}(Y,X)$, with composition reversed. Formally: turn every
arrow around. Note $(\mathcal{C}^{\op})^{\op} = \mathcal{C}$.
\end{definition}

Taking the opposite is trivial \emph{as an operation}; the content of a statement like
``$\Set^{\op} \simeq \CABA$'' is that the formally-reversed category, which has no obvious
description, is equivalent to a category one already knows independently.

\begin{definition}[Functor]
A \emph{functor} $F : \mathcal{C} \to \mathcal{D}$ assigns an object $FX$ to each object $X$
and a morphism $Ff : FX \to FY$ to each $f : X \to Y$, preserving identities and composition:
$F(\mathrm{id}_X) = \mathrm{id}_{FX}$, $F(g \circ f) = Fg \circ Ff$. A functor
$\mathcal{C}^{\op} \to \mathcal{D}$ is called \emph{contravariant} (as a rule on
$\mathcal{C}$ it reverses arrows: $Ff : FY \to FX$, $F(gf) = Ff \circ Fg$).
\end{definition}

\begin{example}[The dual space functor]
$(-)^{*} : \Vect_k^{\op} \to \Vect_k$ sends $V \mapsto \dual{V}$ and $f : V \to W$ to the
\emph{transpose} $\dual{f} : \dual{W} \to \dual{V}$, $\varphi \mapsto \varphi \circ f$. It
reverses arrows because precomposition does.
\end{example}

\begin{example}[The powerset functor]\label{ex:powerset}
$P : \Set^{\op} \to \Set$ sends $X \mapsto 2^{X}$ (the set of subsets) and
$f : X \to Y$ to the \emph{preimage} map $f^{-1} : 2^{Y} \to 2^{X}$. Again arrow-reversing,
because preimage is a pullback operation. Note that the \emph{image} map $f_{*} : 2^X \to 2^Y$
is also functorial but \emph{covariant}, and is not the right functor here: $f^{-1}$ preserves
all unions, intersections and complements, while $f_*$ does not.
\end{example}

\begin{definition}[Natural transformation]
Given functors $F, G : \mathcal{C} \to \mathcal{D}$, a \emph{natural transformation}
$\alpha : F \Rightarrow G$ is a family of morphisms $\alpha_X : FX \to GX$, one per object,
such that for every $f : X \to Y$ the square
\[
\begin{tikzcd}
FX \arrow[r,"\alpha_X"] \arrow[d,"Ff"'] & GX \arrow[d,"Gf"] \\
FY \arrow[r,"\alpha_Y"'] & GY
\end{tikzcd}
\]
commutes. It is a \emph{natural isomorphism} if every $\alpha_X$ is an isomorphism.
\end{definition}

This makes precise the remark in \S1: for the double dual there is a natural transformation
$\Id \Rightarrow (-)^{**}$, given by $v \mapsto \mathrm{ev}_v$ where
$\mathrm{ev}_v(\varphi) = \varphi(v)$, and it is an isomorphism on finite-dimensional spaces.
For the single dual no such transformation can even be written down: $\Id$ and $(-)^{*}$ have
different variance, so the square above does not typecheck.

\begin{definition}[Equivalence, anti-equivalence]
Functors $F : \mathcal{C} \to \mathcal{D}$ and $G : \mathcal{D} \to \mathcal{C}$ form an
\emph{equivalence} $\mathcal{C} \simeq \mathcal{D}$ if $GF \cong \Id_{\mathcal{C}}$ and
$FG \cong \Id_{\mathcal{D}}$ naturally. An \emph{anti-equivalence} or \emph{duality} between
$\mathcal{C}$ and $\mathcal{D}$ is an equivalence $\mathcal{C}^{\op} \simeq \mathcal{D}$.
\end{definition}

Equivalence, not isomorphism, is the correct notion of sameness for categories: it allows
objects to be matched up only up to isomorphism, which is all one ever wants.

\begin{proposition}[Finite-dimensional self-duality]\label{prop:selfdual}
$\FinVect_k^{\op} \simeq \FinVect_k$, via $V \mapsto \dual{V}$, with the natural isomorphism
$V \cong V^{**}$ witnessing the round trip.
\end{proposition}

This is the cheap answer to our question: over any field, finite-dimensional linear algebra is
its own opposite. Nothing new appears. That is suspicious, because in $\Set$ something new
\emph{does} appear, and we should like to know why.

\section{Order, lattices and Boolean algebras, from scratch}\label{sec:lattices}

\begin{definition}[Poset]
A \emph{partially ordered set} (poset) is a set $P$ with a relation $\leq$ that is reflexive
($x \leq x$), antisymmetric ($x \leq y$ and $y \leq x$ imply $x = y$) and transitive.
\end{definition}

\begin{definition}[Meet and join]
In a poset, the \emph{join} $x \vee y$ is the least upper bound of $x$ and $y$ (if it exists)
and the \emph{meet} $x \wedge y$ is the greatest lower bound. A \emph{lattice} is a poset in
which all binary meets and joins exist. It is \emph{bounded} if it has a least element $0$ and
greatest element $1$, and \emph{complete} if \emph{all} subsets (not just finite ones) have
meets and joins.
\end{definition}

\begin{example}\label{ex:lattices}
\begin{enumerate}[nosep,label=(\alph*)]
  \item $2^{X}$, the subsets of a set $X$ ordered by inclusion: $\vee = \cup$, $\wedge = \cap$,
        $0 = \emptyset$, $1 = X$. Complete.
  \item $L(V)$, the subspaces of a vector space ordered by inclusion:
        $W \wedge U = W \cap U$, $W \vee U = W + U$ (\emph{not} $W \cup U$, which is not a
        subspace), $0 = \{0\}$, $1 = V$. Complete.
\end{enumerate}
\end{example}

\begin{definition}[Distributive, modular]\label{def:distmod}
A lattice is \emph{distributive} if
\[
  x \wedge (y \vee z) = (x \wedge y) \vee (x \wedge z) \qquad \text{for all } x,y,z,
\]
and \emph{modular} if the weaker
\[
  x \leq z \;\Longrightarrow\; x \vee (y \wedge z) = (x \vee y) \wedge z
\]
holds. Distributive implies modular.
\end{definition}

\begin{proposition}\label{prop:modular}
$L(V)$ is modular (Dedekind's modular law, valid for submodules of any module), and for
$\dim V \geq 2$ it is \emph{not} distributive.
\end{proposition}

\begin{proof}[Proof of non-distributivity]
Take three pairwise-distinct lines $A, B, C$ in a plane. Then $B \vee C$ is the whole plane,
so $A \wedge (B \vee C) = A$; but $A \wedge B = A \wedge C = 0$, so
$(A \wedge B) \vee (A \wedge C) = 0 \neq A$. Three distinct lines exist in $\FF_q^{2}$ for
every $q \geq 2$, since a plane has $q+1 \geq 3$ lines.
\end{proof}

This single computation is the seed of everything that follows: it is the exact point at which
linear algebra departs from set theory, and it is visibly a statement about the presence of
\emph{more than two} one-dimensional subspaces --- that is, about $q$.

\begin{definition}[Complement, Boolean algebra]
In a bounded lattice, a \emph{complement} of $x$ is an element $y$ with $x \wedge y = 0$ and
$x \vee y = 1$. A lattice is \emph{complemented} if every element has at least one complement.
A \emph{Boolean algebra} is a bounded distributive complemented lattice. In a Boolean algebra
complements are unique, and one writes $\neg x$ or $x^{c}$.
\end{definition}

\begin{warning}\label{warn:complemented}
$L(V)$ \emph{is} complemented --- every subspace has a direct complement --- but the
complement is far from unique, and $W \mapsto W\perpp$ is not in general one of them (that is
exactly the isotropy problem of \S\ref{sec:breaks}). Boolean-ness fails through
\emph{distributivity}, not through the existence of complements.
\end{warning}

\begin{definition}[Atom, atomic, CABA]
In a bounded lattice an \emph{atom} is a minimal nonzero element. A complete lattice is
\emph{atomic} if every element is the join of the atoms below it. $\CABA$ is the category of
\emph{complete atomic Boolean algebras}, with morphisms the maps preserving arbitrary joins,
arbitrary meets and complements.
\end{definition}

\begin{definition}[Boolean homomorphism, $\Bool$, $\FinBool$]
A \emph{Boolean homomorphism} preserves $\vee, \wedge, \neg, 0, 1$. $\Bool$ is the category of
Boolean algebras and homomorphisms; $\FinBool$ the finite ones.
\end{definition}

\begin{definition}[Filter, ultrafilter]
A \emph{filter} on a Boolean algebra $B$ is a nonempty upward-closed subset $F$ closed under
$\wedge$ with $0 \notin F$; it is an \emph{ultrafilter} if for every $x$ exactly one of
$x, \neg x$ lies in $F$. Ultrafilters correspond exactly to homomorphisms $B \to 2$.
\end{definition}

\section{The duality $\Set^{\op} \simeq \CABA$}

\begin{theorem}[Powerset duality]\label{thm:setcaba}
The functor $P : \Set^{\op} \to \CABA$, $X \mapsto 2^{X}$, $f \mapsto f^{-1}$, is an
equivalence of categories. Its inverse sends a complete atomic Boolean algebra to its set of
atoms. Restricting to finite objects, $\FinSet^{\op} \simeq \FinBool$.
\end{theorem}

\begin{proof}[Sketch]
$2^{X}$ is a complete atomic Boolean algebra whose atoms are the singletons, so
$\mathrm{At}(2^{X}) \cong X$ naturally. Conversely, in a complete atomic Boolean algebra $B$
every element is the join of the atoms beneath it, and this decomposition is unique, giving
$B \cong 2^{\mathrm{At}(B)}$ naturally. Morphisms match because a complete Boolean
homomorphism $2^{Y} \to 2^{X}$ is determined by where atoms go, i.e.\ by a function
$X \to Y$.
\end{proof}

\begin{remark}[Two neighbouring theorems]
$\Bool^{\op} \simeq \Stone$, the category of \emph{Stone spaces} (compact, Hausdorff, totally
disconnected topological spaces), by sending $B$ to its space of ultrafilters. This is Stone
duality proper; Theorem~\ref{thm:setcaba} is the ``complete/atomic'' variant in which the
topology becomes invisible. In the finite case all three coincide:
$\FinSet^{\op} \simeq \FinBool$, and finite Stone spaces are just finite sets.
\end{remark}

\section{Why the two dualities differ: dualizing objects}

Both $\FinVect^{\op} \simeq \FinVect$ and $\FinSet^{\op} \simeq \FinBool$ are instances of one
pattern.

\begin{definition}[Dualizing (schizophrenic) object]
Suppose an object $D$ lives in two categories $\mathcal{C}$ and $\mathcal{D}$ ``compatibly''.
Then $\Hom_{\mathcal{C}}(-, D)$ and $\Hom_{\mathcal{D}}(-, D)$ are contravariant functors that
often form an adjunction, and restrict to a duality on the objects for which the canonical map
$X \to \Hom(\Hom(X,D),D)$ is an isomorphism.
\end{definition}

\begin{itemize}
  \item In $\Vect_k$, the dualizing object is $k$: $\dual{V} = \Hom_k(V,k)$. Here $k$ is the
        \emph{unit} of the tensor product and carries no internal structure beyond being the
        ground field, so $\Hom_k(V,k)$ is a bare vector space. The duality stays inside
        $\Vect_k$.
  \item In $\Set$, the dualizing object is $2 = \{\bot,\top\}$: $2^{X} = \Hom_{\Set}(X, 2)$.
        Here $2$ is \emph{not} the unit of the cartesian product (that is the one-point set
        $1$); it is the \emph{subobject classifier}, and it carries an internal Boolean
        algebra structure. So $2^{X}$ inherits that structure pointwise and the duality lands
        \emph{outside} $\Set$, in $\CABA$.
\end{itemize}

\begin{definition}[Subobject]
A morphism $m : U \to X$ is a \emph{monomorphism} (mono) if $mf = mg \Rightarrow f = g$. A
\emph{subobject} of $X$ is an isomorphism class of monos into $X$; $\Sub(X)$ denotes the poset
of subobjects. In $\Set$, $\Sub(X) = 2^{X}$; in $\Vect_k$, $\Sub(V) = L(V)$.
\end{definition}

\begin{definition}[Subobject classifier, topos]
A \emph{subobject classifier} in a category with finite limits is an object $\Omega$ with a
mono $\mathsf{true} : 1 \to \Omega$ such that for every mono $U \rightarrowtail X$ there is a
unique $\chi_U : X \to \Omega$ making
\[
\begin{tikzcd}
U \arrow[r] \arrow[d, rightarrowtail] & 1 \arrow[d,"\mathsf{true}"] \\
X \arrow[r,"\chi_U"'] & \Omega
\end{tikzcd}
\]
a pullback. An \emph{(elementary) topos} is a category with finite limits, exponentials
$Y^{X}$, and a subobject classifier. $\Set$ is a topos with $\Omega = 2$ and $\chi_U$ the
indicator function.
\end{definition}

So in $\Set$ the phrase ``the dualizing object'' and the phrase ``the object that represents
subobjects'' name the \emph{same} object $2$. That coincidence is what makes
$\Set^{\op} \simeq \CABA$ simultaneously a hom-duality and a statement about subobject
lattices. In $\Vect_k$ the coincidence fails --- and it must.

\begin{definition}[Additive and abelian categories]
A category is \emph{additive} if each $\mathcal{C}(X,Y)$ is an abelian group with bilinear
composition, there is a zero object $0$ (both initial and terminal), and finite biproducts
exist. It is \emph{abelian} if additionally every morphism has a kernel and a cokernel, every
mono is a kernel, and every epi is a cokernel. $\Vect_k$ is abelian.
\end{definition}

\begin{definition}[Strict initial object]
An initial object $0$ is \emph{strict} if every morphism $X \to 0$ is an isomorphism. In
$\Set$ the initial object is $\emptyset$, and indeed any function $X \to \emptyset$ forces
$X = \emptyset$. In every topos the initial object is strict.
\end{definition}

\begin{theorem}\label{thm:notopos}
No nontrivial category is both abelian and a topos.
\end{theorem}

\begin{proof}
In an abelian category the initial and terminal objects coincide: $0 \cong 1$. In a topos the
initial object is strict. So for any object $X$, the unique map $X \to 1 \cong 0$ is a
morphism into a strict initial object, hence an isomorphism, hence $X \cong 0$ for every $X$
and the category is equivalent to the one-object one-morphism category.
\end{proof}

\begin{center}
\fbox{\parbox{0.92\textwidth}{\textbf{The structural point.} $\Set$ and $\Vect_q$ sit on
opposite sides of a genuine dichotomy. Subobject lattices in a topos are Heyting algebras
(Boolean, in $\Set$); subobject lattices in an abelian category are modular. By
Theorem~\ref{thm:notopos} no category is both. Hence the $\Set$-style duality and the
$\Vect$-style duality are not two instances of one theorem that one can interpolate naively:
they are two different notions of ``the dual'', which happen to coincide only for $\Set$.}}
\end{center}

\section{Deforming the recipe: subspace lattices and projective geometry}

The \emph{content} of $\Set^{\op} \simeq \CABA$ is: \emph{a set is recovered from its lattice
of subobjects, and that lattice is Boolean.} Let us run that recipe over $\FF_q$.

\begin{definition}[Projective space]
For a vector space $V$ over $k$, the \emph{projective space} $\mathbb{P}(V)$ has as points the
$1$-dimensional subspaces of $V$, as lines the $2$-dimensional subspaces, and so on. When
$\dim V = n$ and $k = \FF_q$ one writes $\PG(n-1, q)$; its \emph{projective dimension} is
$n-1$.
\end{definition}

\begin{definition}[Gaussian binomial coefficient]
For $0 \le k \le n$,
\[
  \qbin{n}{k} \;=\; \prod_{i=0}^{k-1} \frac{q^{\,n-i} - 1}{q^{\,k-i} - 1}.
\]
\end{definition}

\begin{proposition}
The number of $k$-dimensional subspaces of $\FF_q^{n}$ is $\qbin{n}{k}$, and
$\lim_{q \to 1} \qbin{n}{k} = \binom{n}{k}$.
\end{proposition}

\begin{example}
$\qbin{2}{1} = \frac{q^{2}-1}{q-1} = q+1$: a plane over $\FF_q$ has $q+1$ lines, which is why
Proposition~\ref{prop:modular} needed $q \geq 2$. At $q=1$ this is $2 = \binom{2}{1}$.
\end{example}

\begin{theorem}[Fundamental Theorem of Projective Geometry]\label{thm:ftpg}
Let $V, W$ be vector spaces of dimension $\geq 3$ over fields $k, l$. Every isomorphism of
lattices $L(V) \to L(W)$ is induced by a $\sigma$-semilinear isomorphism $V \to W$ for some
field isomorphism $\sigma : k \to l$, and the semilinear map is unique up to a scalar. In
particular
\[
  \Aut\bigl(L(V)\bigr) \;=\; \PGammaL(V) \;=\; \PGL(V) \rtimes \Aut(k).
\]
\end{theorem}

This is the exact analogue of Theorem~\ref{thm:setcaba}: \emph{the object is recovered from
its subobject lattice}. The difference is what kind of lattice one gets.

\begin{center}
\begin{tabular}{lll}
\toprule
Category & $\Sub(X)$ & Recovery theorem \\
\midrule
$\Set$ & Boolean lattice $2^{X}$ & $\Set^{\op} \simeq \CABA$ \\
$\Vect_{\FF_q}$ & projective geometry $L(V) = \PG(n{-}1,q)$ & Fundamental Theorem (\ref{thm:ftpg}) \\
\bottomrule
\end{tabular}
\end{center}

\begin{definition}[Quantum logic]
Birkhoff and von Neumann (1936) proposed reading the subspace lattice of a Hilbert space as a
propositional logic: $\wedge$ = ``and'', $\vee$ = ``or'', $W\perpp$ = ``not''. Because the
lattice is non-distributive, this \emph{quantum logic} differs from classical propositional
logic exactly in the distributive law. The failure computed in
Proposition~\ref{prop:modular} is the same failure.
\end{definition}

\noindent
So: \textbf{the $q$-analogue of $\Bool$ is the category of projective geometries ---
equivalently, finite quantum logic.}

\section{Dualities and polarities: where Part I reappears}

We now identify the anti-automorphisms. In $\Set$ the anti-automorphism of $2^{X}$ is
complementation $S \mapsto S^{c}$. What plays that role for $L(V)$?

\begin{definition}[Correlation, duality, polarity]
An \emph{anti-isomorphism} (or \emph{duality}, or \emph{correlation}) of lattices is an
order-reversing bijection $\delta : L(V) \to L(W)$; equivalently an isomorphism
$L(V)^{\op} \to L(W)$. A \emph{polarity} is a duality $\delta : L(V) \to L(V)$ with
$\delta^{2} = \mathrm{id}$.
\end{definition}

\begin{definition}[Reflexive form]
A sesquilinear form (Definition~\ref{def:sesq}) is \emph{reflexive} if
$\ip{v}{w} = 0 \iff \ip{w}{v} = 0$. Reflexivity is exactly what is needed for
$W \mapsto W\perpp$ to be unambiguous (left and right perp agree).
\end{definition}

\begin{theorem}[Second Fundamental Theorem; Birkhoff--von Neumann, Baer]\label{thm:bvn}
Let $\dim V \geq 3$. Every duality $L(V) \to L(V)$ is of the form $W \mapsto W\perpp$ for a
nondegenerate sesquilinear form, unique up to a scalar. The dualities that are
\emph{polarities} correspond precisely to the \emph{reflexive} forms, and every nondegenerate
reflexive sesquilinear form is, after rescaling by a nonzero scalar, of exactly one of the
following kinds:
\begin{enumerate}[nosep,label=(\roman*)]
  \item symmetric bilinear \quad ($\sigma = \mathrm{id}$, $\ip{w}{v} = \ip{v}{w}$);
  \item alternating bilinear \quad ($\sigma = \mathrm{id}$, $\ip{v}{v} = 0$);
  \item $\sigma$-Hermitian \quad ($\sigma$ an involution $\neq \mathrm{id}$).
\end{enumerate}
(In characteristic $2$ cases (i) and (ii) are not disjoint; see the remark in \S\ref{sec:forms}.)
\end{theorem}

\begin{center}
\fbox{\parbox{0.92\textwidth}{\textbf{Part I and Part II are the same theorem.} The three
families of forms classified in \S\ref{sec:forms} --- symmetric, alternating, Hermitian ---
are \emph{not} a convenient list. By Theorem~\ref{thm:bvn} they are exactly the
implementations of ``op'' for linear algebra: the involutive anti-automorphisms of the
subobject lattice. Asking ``what replaces the inner product over $\FF_q$?'' and asking ``what
is the opposite of $\Vect_{\FF_q}$?'' are the same question, and the classification in
Part~I is the answer to both.}}
\end{center}

\section{Complementation lost}\label{sec:complost}

Here the compressed sentence of \S\ref{sec:breaks} returns one level up.

\begin{definition}[Ortholattice, orthomodular lattice]\label{def:ortho}
An \emph{orthocomplementation} on a bounded lattice is a map $x \mapsto x'$ with
$x'' = x$, \; $x \leq y \Rightarrow y' \leq x'$, \; $x \wedge x' = 0$ and $x \vee x' = 1$. A
lattice with one is an \emph{ortholattice}; it is \emph{orthomodular} if
$x \leq y \Rightarrow y = x \vee (y \wedge x')$. Subspace lattices of Hilbert spaces are
orthomodular; Boolean algebras are exactly the distributive ortholattices.
\end{definition}

Now compare the three regimes for the map $W \mapsto W\perpp$:

\begin{center}
\small
\setlength{\tabcolsep}{4pt}
\begin{tabular}{lccc}
\toprule
 & $\CC$, positive definite & $\FF_q$, nondegenerate & $\FF_1$ \\
\midrule
order-reversing involution & yes & yes & yes \\
$\dim W + \dim W\perpp = \dim V$ & yes & yes & yes \\
$W \wedge W\perpp = 0$ & yes & \textbf{no} & yes \\
so: an orthocomplementation? & yes & \textbf{no} & yes \\
lattice distributive? & no & no & \textbf{yes} \\
resulting logic & orthomodular (quantum) & polarity only & Boolean (classical) \\
\bottomrule
\end{tabular}
\end{center}

Over $\FF_q$ one has \emph{strictly less} than quantum logic: $W \mapsto W\perpp$ is an
order-reversing involution but not an orthocomplementation, because $W \cap W\perpp \neq 0$
whenever $W$ contains an isotropic vector. Boolean complementation is a polarity \emph{plus}
the complement property. Over $\FF_q$ one gets the polarity and loses the complementation ---
which is precisely ``one keeps the duality and loses the splitting'', restated for lattices.

\section{Why it is called quantum logic}\label{sec:qlogic}

The name is Birkhoff and von Neumann's, from \emph{The Logic of Quantum Mechanics}
(\emph{Annals of Mathematics} \textbf{37}, 1936). ``Logic'' means the algebra of
\emph{propositions} about a physical system; ``quantum'' means the propositions that quantum
mechanics permits one to assert. The claim being made is that physics determines which
propositional calculus is correct.

\subsection{Classical propositions form a Boolean algebra}

\begin{definition}[Phase space, classical proposition]
A classical mechanical system has a \emph{phase space} $X$, the set of its possible states
(positions and momenta). A \emph{proposition} about the system is a yes/no question --- ``the
particle lies in region $R$'', ``the energy is below $E$'' --- and is identified with the
subset of $X$ on which it holds.
\end{definition}

Under that identification, ``$P$ and $Q$'' is $P \cap Q$, ``$P$ or $Q$'' is $P \cup Q$, and
``not $P$'' is $P^{c}$. So the propositions of a classical system form a Boolean algebra.
Note that classical logic is not \emph{assumed} here: it is \emph{derived} from the facts
that states are points and propositions are subsets.

\subsection{Quantum propositions are subspaces}

\begin{definition}[Hilbert space]
A \emph{Hilbert space} is a complex vector space $H$ with a positive-definite Hermitian inner
product that is complete in the induced norm. A subspace is \emph{closed} if it contains all
its limits; in finite dimensions every subspace is closed.
\end{definition}

\begin{definition}[Projection, self-adjoint operator]
The \emph{adjoint} $A^{*}$ of a bounded operator is defined by
$\ip{Av}{w} = \ip{v}{A^{*}w}$; $A$ is \emph{self-adjoint} if $A = A^{*}$. A
\emph{(orthogonal) projection} is a self-adjoint $P$ with $P^{2} = P$; projections correspond
bijectively to closed subspaces, $P \leftrightarrow \mathrm{im}\,P$.
\end{definition}

\begin{theorem}[Spectral theorem, projection-valued form]
Every self-adjoint operator $A$ determines a \emph{projection-valued measure} $E_A$: to each
Borel set $\Delta \subseteq \RR$ it assigns a projection $E_A(\Delta)$, additively on disjoint
unions, with $E_A(\RR) = \mathrm{id}$.
\end{theorem}

The projection $E_A(\Delta)$ \emph{is} the proposition ``the value of the observable $A$ lies
in $\Delta$''. Hence:
\begin{center}
\fbox{\parbox{0.9\textwidth}{Quantum propositions are the closed subspaces of $H$, and
\[
  \text{``and''} = W \cap U, \qquad
  \text{``or''} = \overline{W + U}, \qquad
  \text{``not''} = W\perpp .
\]}}
\end{center}
A state $\psi$ makes the proposition $W$ \emph{definitely true} if $\psi \in W$,
\emph{definitely false} if $\psi \in W\perpp$, and otherwise assigns it only the probability
$\lVert P_W \psi\rVert^{2}$ (the \emph{Born rule}).

The decisive entry is the middle one. ``Or'' is \emph{not} union, because a union of subspaces
is not a subspace; one must take the closed span, and the span contains vectors lying in
neither disjunct. Those are the superpositions.

\subsection{The failure of distributivity, read physically}

Take a qubit, $H = \CC^{2}$, with the usual notation $\lvert 0\rangle, \lvert 1\rangle$ for
the $z$-spin eigenstates and $\lvert +\rangle = (\lvert 0\rangle + \lvert 1\rangle)/\sqrt{2}$
for an $x$-spin eigenstate. Put
\[
  a = \text{``spin is } {+}x\text{''} = \langle \lvert +\rangle\rangle, \quad
  b = \text{``spin is } {+}z\text{''} = \langle \lvert 0\rangle\rangle, \quad
  c = \text{``spin is } {-}z\text{''} = \langle \lvert 1\rangle\rangle .
\]
Then $b \vee c = H$, the tautology ``the $z$-spin is up or down'', so
\[
  a \wedge (b \vee c) = a \wedge 1 = a,
  \qquad\text{but}\qquad
  (a \wedge b) \vee (a \wedge c) = 0 \vee 0 = 0 .
\]
In words: ``the spin is $+x$, \emph{and} the $z$-spin is either up or down'' is true, since
the second conjunct is a tautology; whereas ``either (the spin is $+x$ and $z$-up) or (the
spin is $+x$ and $z$-down)'' is a disjunction of two contradictions, there being no
simultaneous eigenstate of $x$- and $z$-spin. Distributing across the disjunction converts a
truth into a falsehood.

The physical content is: \emph{a disjunction can be true without either disjunct being true}.
$\lvert +\rangle$ lies in the span of $\{\lvert 0\rangle, \lvert 1\rangle\}$ but in neither.
Superposition, stated as a fact about a lattice.

\begin{remark}[This is Proposition~\ref{prop:modular}]\label{rem:samecomputation}
The computation just performed is the one proving $L(V)$ non-distributive: three
pairwise-distinct lines $A, B, C$ in a plane satisfy $A \wedge (B \vee C) = A$ while
$(A \wedge B) \vee (A \wedge C) = 0$. Here the plane is $\CC^{2}$ and the three lines are
$\lvert 0\rangle$, $\lvert 1\rangle$, $\lvert +\rangle$.

Over $\FF_2$ a plane has exactly three lines, spanned by $(1,0)$, $(0,1)$ and $(1,1)$ ---
that is, by $\lvert 0\rangle$, $\lvert 1\rangle$ and $\lvert 0\rangle + \lvert 1\rangle$. So
the smallest non-distributive quantum logic is the $\FF_2$ qubit, and the hypothesis
``$q \geq 2$'' in the proof of Proposition~\ref{prop:modular} is exactly the hypothesis that
superposition exists at all.
\end{remark}

\subsection{Is it really \emph{logic}? The case against}

As a logic it is impoverished, and this should be said plainly.
\begin{itemize}[nosep]
  \item There is no satisfactory conditional. The deduction theorem fails, and the candidate
        implications (the Sasaki hook $a \to b := a' \vee (a \wedge b)$ and its rivals) each
        fail some property one wants.
  \item Orthomodular lattices have no well-behaved tensor product, so composing two systems
        --- the very place quantum mechanics is most distinctive --- is awkward.
  \item A deflationary reading is available: nothing about \emph{inference} has changed; what
        changed is which propositions are simultaneously meaningful. On that reading the
        subject is projective geometry with a suggestive vocabulary.
\end{itemize}
Putnam's \emph{Is Logic Empirical?} (1968) argued for the strong reading, that quantum
mechanics revises logic as relativity revised geometry. He later retracted it.

\subsection{The case for: the lattice determines the physics}

The name is nonetheless earned, because the propositional lattice turns out to contain the
entire theory --- a startling amount of content for a partially ordered set.

\begin{theorem}[Gleason, 1957]
Let $\dim H \geq 3$. Every countably additive probability measure on the lattice of
projections of $H$ has the form $P \mapsto \mathrm{tr}(\rho P)$ for a unique density operator
$\rho$.
\end{theorem}

So the Born rule need not be postulated: it is forced by the lattice together with countable
additivity. (The hypothesis $\dim H \geq 3$ is essential --- the qubit used above admits exotic
measures --- and it is the same dimension bound as in the Fundamental
Theorem~\ref{thm:ftpg}.)

\begin{definition}[Division ring, quaternions]
A \emph{division ring} satisfies all the field axioms except commutativity of multiplication.
The \emph{quaternions} $\HH$ form the four-dimensional real division ring with basis
$1, i, j, k$ and $i^{2} = j^{2} = k^{2} = ijk = -1$, with conjugation
$\overline{a + bi + cj + dk} = a - bi - cj - dk$ as involution.
\end{definition}

\begin{theorem}[Piron 1964; Sol\`er 1995]\label{thm:pironsoler}
A complete, atomistic, irreducible orthomodular lattice with the covering property and
dimension $\geq 4$ is isomorphic to the lattice of closed subspaces of a Hermitian space over
some division ring with involution. If it moreover admits an infinite orthonormal sequence,
that division ring is $\RR$, $\CC$ or $\HH$ with its standard involution, and the space is a
Hilbert space.
\end{theorem}

Piron's step is, in essence, the Fundamental Theorem of Projective Geometry
(Theorem~\ref{thm:ftpg}) together with the classification of polarities
(Theorem~\ref{thm:bvn}). Which means the three families of Part~I reappear a third time ---
now as the complete list of candidate quantum mechanics:

\begin{center}
\small
\setlength{\tabcolsep}{5pt}
\begin{tabular}{lll}
\toprule
polarity type (Thm.~\ref{thm:bvn}) & division ring, involution & physics \\
\midrule
symmetric bilinear & $\RR$, trivial involution & real quantum mechanics \\
Hermitian & $\CC$, complex conjugation & standard quantum mechanics \\
Hermitian & $\HH$, quaternionic conjugation & quaternionic quantum mechanics \\
\textbf{alternating} & --- & \textbf{excluded} \\
\bottomrule
\end{tabular}
\end{center}

\begin{proposition}[Why symplectic quantum mechanics is impossible]
An alternating form satisfies $\ip{v}{v} = 0$ for every $v$, so every subspace meets its own
perp: $W \wedge W\perpp \neq 0$ for all $W \neq 0$. Hence $W \mapsto W\perpp$ is never an
orthocomplementation (Definition~\ref{def:ortho}), and the lattice is never an ortholattice.
\end{proposition}

\begin{center}
\fbox{\parbox{0.92\textwidth}{\textbf{Positive-definiteness is what selects two of the three
families and discards the symplectic one.} The classification of \S\ref{sec:forms} is
simultaneously: (i) the answer to ``what replaces the inner product over $\FF_q$'', (ii) by
Theorem~\ref{thm:bvn}, the list of implementations of ``op'' for linear algebra, and (iii) by
Theorem~\ref{thm:pironsoler}, the list of possible quantum mechanics. One classification,
three readings.}}
\end{center}

\subsection{Over $\FF_q$: the logic without the statistics}

What does the finite-field case inherit? Exactly the lattice, and nothing measure-theoretic.
The obstruction is Remark~\ref{rem:ordering}: no ordering, hence no probabilities.

\begin{center}
\begin{tabular}{ll}
\toprule
kept over $\FF_q$ & lost over $\FF_q$ \\
\midrule
superposition (Remark~\ref{rem:samecomputation}) & the Born rule: $\lVert P_W\psi\rVert^{2}$
  needs $[0,1] \subset \RR$ \\
non-distributivity & Gleason's theorem: no measures to classify \\
entanglement, no-cloning, contextuality & orthomodularity itself (table in
  \S\ref{sec:complost}) \\
\bottomrule
\end{tabular}
\end{center}

Amplitudes exist over $\FF_q$; magnitudes do not. This is a live research area ---
\emph{modal quantum theory} (Schumacher and Westmoreland), quantum mechanics over finite
fields, in which the only predicates are \emph{possible} and \emph{impossible} rather than
probabilities. Remarkably, analogues of teleportation, superdense coding, no-cloning and a
Bell-type contextuality theorem all survive there. It is a clean experiment separating the
structural features of quantum mechanics from those requiring probability.

\subsection{A wrinkle: von Neumann's own retreat}

Von Neumann grew dissatisfied with the projection lattice of $B(H)$. For infinite-dimensional
$H$ that lattice is orthomodular but \emph{not modular} (Definition~\ref{def:distmod}), and
modularity is what supports a well-behaved dimension function. Wanting a dimension taking
values continuously in $[0,1]$ --- a canonical \emph{a priori} probability --- he was led to
type $\mathrm{II}_1$ factors and continuous geometry, and away from quantum logic toward what
are now called von Neumann algebras.

The wrinkle for us: over $\FF_q$ every space is finite-dimensional, so $L(V)$ \emph{is}
modular (Proposition~\ref{prop:modular}). The finite-field case sits on von Neumann's
preferred side of that split. It has the modularity he wanted, and lacks --- for the entirely
separate reason that $\FF_q$ is unordered --- the measure he wanted it for.


\newpage
\part{The limit $q \to 1$}

\section{The field with one element}

\begin{warning}
There is no field with one element: a field requires $1 \neq 0$. ``$\FF_1$'' names a
programme, not an object --- the search for a base beneath $\Spec \ZZ$ over which
$q$-analogue formulas make sense at $q = 1$. Implementations include Tits' original
observation (1957), Deitmar's monoid schemes, Kapranov--Smirnov, Connes--Consani, Borger's
$\Lambda$-rings and To\"en--Vaqui\'e. They agree on the combinatorics below and differ
elsewhere.
\end{warning}

\begin{definition}[Pointed set; the standard model]
A \emph{pointed set} is a pair $(X, \ast)$ with $\ast \in X$; a morphism of pointed sets
preserves the basepoint. $\FinSet_{\ast}$ denotes the category of finite pointed sets. In the
Deitmar / Kapranov--Smirnov model, an \emph{$\FF_1$-vector space} is a finite pointed set (the
basepoint is $0$), and an \emph{$\FF_1$-linear map} is a pointed map. There is no addition ---
$\FF_1$ has no additive structure to inherit --- and scalar multiplication is trivial because
$\FF_1^{\times} = \{1\}$.
\end{definition}

\begin{proposition}[$\GL$ at $q = 1$]
$\GL_n(\FF_1) = S_n$, the symmetric group: the $\FF_1$-linear automorphisms of a pointed set
with $n$ non-basepoint elements are the permutations of those elements. Numerically,
\[
  |\GL_n(\FF_q)| \;=\; \prod_{i=0}^{n-1}\bigl(q^{n} - q^{i}\bigr)
  \;=\; q^{\binom{n}{2}} (q-1)^{n} \,[n]_q!,
  \qquad [n]_q! = \prod_{i=1}^{n} \frac{q^{i}-1}{q-1},
\]
so after dividing out the $n$ copies of the scalar group $\FF_q^{\times}$ --- which
degenerates to the trivial group --- and setting $q = 1$ one gets $n!$.
\end{proposition}

\begin{remark}[Tits' observation]
The building of $\GL_n(\FF_q)$ is the flag complex of the subspace lattice $L(\FF_q^{n})$; at
$q=1$ it degenerates to the Coxeter complex of $S_n$, i.e.\ to a single apartment. ``The
building becomes the apartment'' and ``the subspace lattice becomes the Boolean lattice'' are
the same statement: a frame (unordered basis up to scalars) is exactly a set.
\end{remark}

\section{Isotropy vanishes, and Boolean logic reappears}

\begin{proposition}[Counts of nonzero isotropic vectors]\label{prop:isocount}
For a nondegenerate form on a space over $\FF_q$, the number of nonzero isotropic vectors is:
\begin{center}
\begin{tabular}{lll}
\toprule
form & count & value at $q=1$ \\
\midrule
symplectic $\Sp_{2m}(q)$ & $q^{2m}-1$ & $0$ \\
orthogonal $O_{2m}^{+}(q)$ & $(q^{m}-1)(q^{m-1}+1)$ & $0$ \\
orthogonal $O_{2m}^{-}(q)$ & $(q^{m}+1)(q^{m-1}-1)$ & $0$ \\
orthogonal $O_{2m+1}(q)$ & $q^{2m}-1$ & $0$ \\
unitary $U_{n}(q)$ & $\bigl(q^{n}-(-1)^{n}\bigr)\bigl(q^{n-1}-(-1)^{n-1}\bigr)$ & $0$ \\
\bottomrule
\end{tabular}
\end{center}
Every one of them vanishes at $q = 1$.
\end{proposition}

\begin{proof}[Why this must happen]
Isotropy requires \emph{addition}. To have $\ip{v}{v} = 0$ with $v \neq 0$ one needs a sum of
nonzero terms to cancel, and $\FF_1$ has no addition at all. The vanishing of the counts is
the arithmetic shadow of that structural absence.
\end{proof}

\begin{corollary}[The chain closes]
At $q = 1$:
\begin{enumerate}[nosep,label=(\alph*)]
  \item there are no isotropic vectors, so $W \wedge W\perpp = 0$ and the polarity
        $W \mapsto W\perpp$ becomes a genuine complementation;
  \item the subspace lattice $L_n(q)$ becomes the Boolean lattice $2^{[n]}$, hence
        distributive;
  \item distributive $+$ complemented $=$ Boolean.
\end{enumerate}
Therefore the $q$-analogue of $\Bool$ degenerates to $\Bool$, and
$\Set^{\op} \simeq \CABA$ is recovered:
\[
  \Vect_{\FF_q} \;\xrightarrow{\ \text{subobject lattices}\ }\;
  \text{projective geometry / quantum logic} \;\xrightarrow{\ q \to 1\ }\; \Bool .
\]
\end{corollary}

\subsection*{The $q \to 1$ dictionary}

\begin{center}
\small
\setlength{\tabcolsep}{4pt}
\begin{tabular}{lll}
\toprule
& at $q$ & at $q = 1$ \\
\midrule
object & $\FF_q$-vector space $\FF_q^{n}$ & pointed set, $n$ non-base points \\
number of $k$-subobjects & $\qbin{n}{k}$ & $\binom{n}{k}$ \\
subobject lattice & $L_n(q)$: modular, not distributive & $2^{[n]}$: Boolean \\
points of $\PG(n{-}1,q)$ & $1 + q + \cdots + q^{n-1}$ & $n$ \\
automorphisms of the lattice & $\PGammaL_n(q)$ & $S_n$ \\
nonzero isotropic vectors & see Prop.~\ref{prop:isocount} & $0$ \\
$W \mapsto W\perpp$ & polarity, not a complementation & Boolean complementation \\
building & Tits building of $\GL_n(\FF_q)$ & Coxeter complex of $S_n$ \\
\bottomrule
\end{tabular}
\end{center}

\section{An aside: the prime at infinity is not $\FF_1$}

\begin{definition}[Absolute value, place]
An \emph{absolute value} on a field $k$ is $|\cdot| : k \to \RR_{\geq 0}$ with $|x| = 0 \iff
x = 0$, $|xy| = |x||y|$ and $|x+y| \leq |x| + |y|$. Two are \emph{equivalent} if one is a
positive power of the other; an equivalence class is a \emph{place}. The \emph{$p$-adic}
absolute value is $|p^{a} u/v|_p = p^{-a}$ for $u,v$ coprime to $p$.
\end{definition}

\begin{theorem}[Ostrowski]
Every nontrivial place of $\QQ$ is either $|\cdot|_p$ for a prime $p$, or the usual
archimedean $|\cdot|_{\infty}$. The completions are $\QQ_p$ and $\RR$, and the product formula
$\prod_{v} |x|_v = 1$ holds for $x \in \QQ^{\times}$.
\end{theorem}

The place $\infty$ is thus a genuine place of $\QQ$, ``the prime at infinity'', with completion
$\RR$ and residue characteristic $0$. It is an \emph{archimedean completion}. $\FF_1$, by
contrast, is a proposed \emph{degeneration}: a base such that $\Spec \ZZ \to \Spec \FF_1$ is
terminal, so that the fibre product $\Spec\ZZ \times_{\Spec \FF_1} \Spec\ZZ$ exists and a
Weil-style proof of the Riemann Hypothesis might be imitated (Deninger, Manin). These are
different directions: one completes, the other degenerates.

They are nonetheless productively confused, because both are ``the extra point one bolts onto
$\Spec \ZZ$'', and because Connes--Consani genuinely knit them together via the ad\`ele class
space, where the archimedean place and the $\FF_1$-structure appear in one construction. The
confusion is fertile; it is still a confusion.

\section{Three answers, preserved}

The phrase ``the opposite of $\Set$ is $\Bool$'' bundles two claims which coincide only because
$\Set$ is a topos. Unbundled, the $q$-deformation goes three ways, and I do not think they
reconcile.

\begin{description}[leftmargin=1.8em,style=nextline]
\item[(1) The literal answer: linear algebra is self-dual.]
$\FinVect_q^{\op} \simeq \FinVect_q$ (Proposition~\ref{prop:selfdual}). For infinite
dimensions one needs topology:

\begin{definition}[Linearly compact vector space]
A topological $k$-vector space ($k$ discrete) is \emph{linearly topologised} if $0$ has a
neighbourhood basis of open subspaces, and \emph{linearly compact} if in addition it is
Hausdorff, complete, and every family of closed cosets with the finite-intersection property
has nonempty intersection. Equivalently: it is an inverse limit of finite-dimensional discrete
spaces.
\end{definition}

\begin{theorem}[Lefschetz duality]
$\Vect_k^{\op} \simeq \LinComp_k$, via $V \mapsto \dual{V}$ with the topology of pointwise
convergence.
\end{theorem}

\begin{definition}[$\Ind$ and $\Pro$]
$\Ind(\mathcal{C})$ is the category of formal filtered colimits of objects of $\mathcal{C}$;
$\Pro(\mathcal{C}) = \Ind(\mathcal{C}^{\op})^{\op}$, formal cofiltered limits.
\end{definition}

Since $\Vect_k = \Ind(\FinVect_k)$ we get $\Vect_k^{\op} = \Pro(\FinVect_k^{\op}) =
\Pro(\FinVect_k) = \LinComp_k$; and since $\Set = \Ind(\FinSet)$ we get
$\Set^{\op} = \Pro(\FinBool) \simeq \CABA$, while $\Bool^{\op} = \Pro(\FinSet) = \Stone$. All
three dualities are the same formal move --- pass to the pro-completion --- and in that sense
$\LinComp_{\FF_q}$ \emph{is} the $q$-analogue of $\CABA$.

\item[(2) The structural answer: projective geometry.]
Deform the recipe rather than the statement. Subobject lattices give
$\Set \rightsquigarrow \Bool$ and $\Vect_q \rightsquigarrow \PG(n-1,q)$; recovery is
Theorem~\ref{thm:ftpg} in place of Theorem~\ref{thm:setcaba}; the anti-automorphisms are
polarities, classified by Theorem~\ref{thm:bvn} into exactly the three families of Part~I;
and the whole thing degenerates to $\Bool$ at $q = 1$. This is the answer that makes the
question's chain close.

\item[(3) The answer at $q = 1$ itself.]
The opposite of $\FF_1$-linear algebra is $\FinSet_{\ast}^{\op}$, which is Segal's category
$\Gamma$ --- the indexing category for $\Gamma$-spaces and hence for $E_{\infty}$ (coherently
commutative) structures in homotopy theory. Concretely:
\begin{proposition}
$\FinSet_{\ast}^{\op} \simeq \FinBool_{/2}$, the category of finite Boolean algebras equipped
with a homomorphism to $2$, i.e.\ with a distinguished ultrafilter.
\end{proposition}
\begin{proof}
$\FinSet_{\ast} = \FinSet_{1/}$, the coslice under the one-point set. Opposites exchange
coslice and slice: $(\mathcal{C}_{1/})^{\op} = (\mathcal{C}^{\op})_{/1}$. Under
$\FinSet^{\op} \simeq \FinBool$ the object $1$ corresponds to $P(1) = 2$. Finally a
homomorphism $B \to 2$ is the same as an ultrafilter on $B$.
\end{proof}
So even the basepoint --- the $0$ of $\FF_1$-linear algebra --- has a Boolean shadow.
\end{description}

\begin{center}
\fbox{\parbox{0.92\textwidth}{\textbf{Mai nafka minah.} Answer (1) is right if one asks
\emph{what has the same shape}; answer (2) is right if one asks \emph{what does the same job};
answer (3) is right if one asks \emph{what happens at the degenerate point}. They diverge
because $\Set$ conflates its dualizing object with its subobject classifier and $\Vect_q$
cannot, by Theorem~\ref{thm:notopos}. The disagreement should be kept, not resolved.}}
\end{center}

\section*{Coda}

\begin{center}
\begin{tabular}{ccc}
\toprule
$q$ general & $\longrightarrow$ & $q = 1$ \\
\midrule
$\Vect_{\FF_q}$ && $\FinSet_{\ast}$ \\
projective geometry $\PG(n{-}1,q)$ && Boolean algebra $2^{[n]}$ \\
polarity $W \mapsto W\perpp$ && complementation $S \mapsto S^{c}$ \\
symmetric / alternating / Hermitian && (isotropy empty) \\
quantum logic && classical logic \\
\bottomrule
\end{tabular}
\end{center}

\bigskip
\noindent
\emph{One keeps the duality and loses the splitting; at $q=1$ the splitting comes back, and
with it, Boole.}

\end{document}
